\section{Limitations and conclusion}
\label{sec:conclusion}
The questions are model-drafted, forced-choice diagnostics rather than independent user requests. The selected source review found invalid and insufficient items, especially in the reviewed CUAD and DeFi cases; its counts cannot estimate defect prevalence. Removing those found items barely changes the reported AUROCs, but unreviewed keys remain a threat. $R$ also predicts errors on usual-meaning and control terms, so it cannot identify an error's semantic cause by itself. The definition-allocation experiments are too limited to show an advantage over frequency, and the historical retrieval comparison is not a held-out request test.

The supported implication is a division of labour for systems that use local documents. At corpus preparation time, source-linked diagnostic questions can rank uses for human checking without asking every deployed reader to discover its own missing context. The checked definition should remain attached to its source and scope, because its wording changes both repairs and harms. At request time, retrieval still has to bring in the relevant definition and any linked clauses; the present score is not a substitute for that step. Evaluations should test missing local evidence as well as conflicts with evidence already shown. This is a useful offline diagnostic and a concrete design requirement for a semantic layer, while the benefit of an automatic definition-selection policy on new requests remains to be demonstrated.
\label{main:end}

